% !TEX root = ../../tree_notes.tex
% =====================================================================
%  2.4  Structural Transformation                       sec:c02-transform
%  Exports: def:c02-return-newroot  def:c02-flip  def:c02-mutate-rebuild
%           thm:c02-invert-invol  thm:c02-rewrite-frame
%           lst:c02-transform-skeleton  ex:c02-invert-trace
%           fig:c02-invert  fig:c02-where
%           prob:c02-print  prob:c02-create-tree  prob:c02-flip-eq
%           prob:c02-delete-forest
% =====================================================================
\section{Structural Transformation: Rewiring the Tree}\label{sec:c02-transform}

\begin{previouslybox}
\begin{enumerate}[nosep, leftmargin=1.6em]
  \item In \probref{c02-prune} a \emph{boolean} return value triggered a pointer
        write. What exactly did that value signal, and why did the write have to
        happen \emph{after} the recursive calls? (\secref{c02-bottomup})
  \item On $R$, what are $\height(2)$ and $\height(3)$ (in edges), and which
        channel -- up or down -- did that number travel? (\secref{c02-bottomup})
  \item \probref{c02-path-sum-iii} had to \emph{undo} two hashmap entries on the
        way back up. What general rule about mutating shared state inside a
        recursion does that teach you -- and where will it bite again today?
        (\secref{c02-topdown})
\end{enumerate}
\end{previouslybox}

\begin{topicintro}
For three sections the tree has been a read-only object: you carried numbers down
and returned numbers up, but the pointers never moved. This section points the
recursion \emph{at the tree itself}. A \term{structural transformation} changes
which node is whose child -- swap two children, cut a subtree, graft a new node,
rebuild the whole shape from a description. The unifying idiom is small but sharp:
\textbf{make the recursive call \emph{return the new root} of the subtree it
transformed}, so the parent can decide what to do with it. Master that one line
and Invert, Prune, Delete-and-Split, and the construct-from-description family
all become the same twelve lines with a different rewrite.

\textbf{Prerequisites.} The return contract of \secref{c02-bottomup} (a call hands
a value \emph{up}); the base-case discipline of \secref{c02-model}; the
return-that-rewires preview you already met in \probref{c02-prune}.

\textbf{After this section you can}
\begin{itemize}[nosep]
  \item write the ``return the transformed subtree root'' skeleton and use it to
        swap, prune, and splice pointers without leaking the old shape;      \hfill\lv{3}
  \item place the rewrite on the correct side of the recursion -- act on the way
        \emph{down} when the parent decides, on the way \emph{up} when the
        children's results decide;                                            \hfill\lv{3}--\lv{4}
  \item rebuild a tree from indirect descriptions (layouts, parent/child rules)
        by treating construction as recursion with a returned root;           \hfill\lv{3}--\lv{4}
  \item recognise an \emph{equivalence-under-flip} question (``can flips make
        these trees equal?'') and prove the greedy/recursive test correct.     \hfill\lv{4}
\end{itemize}
\end{topicintro}

\begin{recognitionbox}[Pattern card -- structural transformation]
\reckey{``invert / mirror / flip the tree'', ``prune / remove / delete nodes'',
``return the root after\ldots'', ``rebuild the tree from\ldots'', ``split into
forests'', ``can swapping children make them equal?'' -- the \emph{shape}, not a
number, is the output.}
\recuse{
  \item the answer is itself a tree (or its root pointer): invert, prune,
        delete-and-return-forest -- return the new subtree root (\defref{c02-return-newroot});
  \item the tree must be \emph{built} from a specification rather than given
        (a grid of descriptions, two traversals): construct top-down, hand the
        assembled root back;
  \item a \emph{yes/no} question whose witness is a set of flips: compare subtree
        invariants that a flip cannot change (\probref{c02-flip-eq}).
}
\recnotuse{
  \item you only need a \emph{number} about the tree (sum, height, count) and the
        pointers stay put -- that is pure bottom-up, \secref{c02-bottomup};
  \item the task is a fixed \emph{visit order} (print inorder/level-order) with no
        pointer change -- \chref{c03};
  \item children are kept \emph{value-ordered} and must stay so for search --
        that is a BST operation, \chref{c05}, with balance invariants this
        chapter never touches.
}
One-line test: \emph{does a correct program leave the \emph{pointer graph}
different from how it started? If yes -- swap/cut/add edges -- this is a
transformation and every call should return the root of the subtree it rebuilt.}
\end{recognitionbox}

\subsection{Intuition}

Recall Invert Binary Tree (LC 226), which you solved in Chapter~1: swap the two
children of \emph{every} node. Its whole program is three lines, and the reason it
is three lines is that the recursion returns the node it just finished flipping,
so the caller can reattach it. Read $R$ and its mirror in \figref{c02-invert}: the
\emph{shape} changes, the values stay welded to their nodes.

\begin{figure}[H]
\centering
\begin{tabular}{@{}c c c@{}}
\begin{forest} snaptree
[1 [2 [4] [5]] [3 [,phantom] [6]] ]
\end{forest}
& $\xrightarrow{\ \text{invert}\ }$ &
\begin{forest} snaptree
[1 [3 [6] [,phantom]] [2 [5] [4]] ]
\end{forest} \\
{\footnotesize original $R$} & & {\footnotesize every node's children swapped}
\end{tabular}
\caption{Inverting $R$. Node $3$ had a lone \emph{right} child $6$; after the
flip $6$ hangs on the \emph{left}. Node $2$'s pair $\{4,5\}$ becomes $\{5,4\}$.
The values never move -- only the two pointers at each node change.}\label{fig:c02-invert}
\end{figure}

Now the key idea that generalises. A transformation recursion has the same
topology as the bottom-up recursions of \secref{c02-bottomup}, but the returned
value is a \emph{pointer}, not a number:

\begin{center}
\begin{tikzpicture}[diagram, node distance=3mm and 14mm]
  \node[boxnode, fill=softpurple!10, draw=softpurple] (s)
    {\code{dfs(v)}\\returns the\\new $\subtree(v)$ root};
  \node[annot, below=1mm of s.south, align=center, font=\scriptsize]
    {parent does \emph{whatever the rewrite is}\\using the returned child roots,\\then returns \emph{itself} (or a stand-in)};
\end{tikzpicture}
\end{center}

Because the return type is a node pointer, the recursion can do things a numeric
return never could: \emph{delete} a subtree (return \code{nullptr}), \emph{split}
a tree into several roots (\probref{c02-delete-forest}), or \emph{assemble} nodes
that did not exist before (\probref{c02-create-tree}). Everything is the same
skeleton -- only the rewrite line and the return change.

\subsection{Definitions and Notation}

\begin{definition}{Flip at a node}{c02-flip}
A \term{flip} (or \term{mirror step}) at node $v$ is the assignment
$v\to\mathit{left} \leftrightarrow v\to\mathit{right}$: the two child pointers are
exchanged, nothing else changes. A \term{flip set} is a choice, for some nodes, of
whether to flip them; flipping a node's descendants is independent of flipping the
node.
\end{definition}

\noindent\textbf{In plain words.} A flip is a swap of two pointers, never a move of
values. One flip is local; ``invert the tree'' means flip at \emph{every} node.
Because each node is flipped at most once and children are separate objects, the
order in which you flip across nodes does not matter.

\begin{definition}{Return-the-new-root contract}{c02-return-newroot}
In a \term{transforming recursion} the call $\mathit{dfs}(v)$ returns the pointer
that should stand in $v$'s former place in $v\!\to\!\mathit{parent}$'s child slot:
often $v$ itself (possibly rewired), sometimes \code{nullptr} (delete), sometimes a
different node (splice/replace). The caller \emph{assigns} the returned pointer into
the child slot it recursed from.
\end{definition}

\noindent\textbf{In plain words.} ``Whatever you turn this subtree into, hand its
top back to me and I will plug it in where the old one was.'' Deletion is the
special case ``return nothing''; a forest split is ``return nothing \emph{and}
report a new tree here''. This one sentence collapses \probref{c02-delete-forest}
into three lines.

\begin{definition}{Mutate in place vs.\ rebuild}{c02-mutate-rebuild}
A transformation \term{mutates in place} when it edits the pointers of the given
nodes (invert, prune). It \term{rebuilds} when it allocates a fresh tree from a
specification, leaving the input intact (construct-from-traversals, \probref{c02-create-tree}).
Both use a returning recursion; they differ only in whether the returned node is
an \emph{existing} one or a \emph{newly allocated} one.
\end{definition}

\noindent\textbf{In plain words.} Ask ``am I allowed to trash the input?'' If yes,
edit pointers (cheaper, and LeetCode hands you ownership). If the input must
survive -- or must be assembled from scratch -- you build new nodes. The recursion
shape is identical either way.

\begin{notebox}[Mutate in place vs.\ rebuild]
\textbf{Mutate in place} edits the pointers of the given nodes and hands back the
same node (or \code{nullptr} to delete it): invert, prune, delete-and-split.
\textbf{Rebuild fresh} allocates new nodes from a specification, leaving the
read-only input intact, and returns a brand-new root: \probref{c02-create-tree},
construct-from-traversals. Both run on the same returning recursion
(\defref{c02-return-newroot}); they differ only in whether the returned node
already existed.
\end{notebox}

\subsection{Core Properties}

\begin{theorem}{Inversion is an Involution}{c02-invert-invol}
Let $\mathcal{I}(T)$ be the tree obtained by flipping every node of $T$. Then
$\mathcal{I}(\mathcal{I}(T)) = T$ for every binary tree $T$; equivalently,
$\mathcal{I}$ is its own inverse.
\end{theorem}
\begin{proofbox}
Structural induction on $\size$. Base $T=\nil$: $\mathcal{I}(\nil)=\nil$, trivially
equal. Step $T=(v,L,R)$: by definition $\mathcal{I}(T)$ has root $v$ with children
$\mathcal{I}(R)$ (left) and $\mathcal{I}(L)$ (right). Applying $\mathcal{I}$ again
swaps them back and, by the inductive hypothesis, $\mathcal{I}(\mathcal{I}(L))=L$
and $\mathcal{I}(\mathcal{I}(R))=R$, so $\mathcal{I}(\mathcal{I}(T))$ has children
$(L,R)$ -- exactly $T$. A double flip restores each pair of pointers, and the
operations on distinct nodes commute, so order of visiting is irrelevant.
\end{proofbox}
\begin{insightbox}
The proof uses only ``swap is self-inverse'' and ``children are independent'' --
which is why invert needs no bookkeeping: no flag to remember, no restore step
(\emph{cf.}~the hashmap you had to undo in \probref{c02-path-sum-iii}). Any transformation
whose per-node action is self-inverse and local inherits this safety.
\end{insightbox}

\begin{theorem}{Where the Rewrite Goes}{c02-rewrite-frame}
Let a transformation act independently at each node with an action that depends
on (i) the node's own data and (ii) the \emph{already-transformed} subtrees of its
children. Then performing the action \emph{after} the recursive calls
(post-order) is correct. Dually, if the action at $v$ depends on a decision made
by $v$'s \emph{parent} (e.g.\ ``flip here only if the parent told you to''), that
decision must be carried \emph{into} the call and the action taken on the way
\emph{down} (pre-order), before recursing.
\end{theorem}
\begin{proofbox}
Post-order case: the action at $v$ reads the new child roots; recursing first
guarantees, by \thmref{c02-postorder} (structural induction), those children are
correct, so $v$'s action has valid inputs and returns a correct $\subtree(v)$.
Pre-order case: the parent's decision is a function of the path to $v$, not of
$\subtree(v)$'s interior, so it is fully known on entry; applying it before
recursing lets each child inherit the state its own action needs. Placing a
child-dependent rewrite in pre-order would read \emph{stale} (untransformed)
children, and placing a parent-decided flip in post-order would flip after the
children already committed to the wrong orientation -- both give a wrong shape,
verifiable on a two-level tree.
\end{proofbox}

\subsection{The Transform Skeleton}

\begin{lstlisting}[caption={The transforming recursion: SLOT 1 (recurse), SLOT 2 (the rewrite), SLOT 3 (what to return). Which side SLOT 2 goes on is decided by \thmref{c02-rewrite-frame}.}, label={lst:c02-transform-skeleton}]
TreeNode* dfs(TreeNode* v /*, carried state if pre-order */) {
    if (!v) return nullptr;               // SLOT 0: null in, null out
    // [PRE-ORDER rewrite here if the PARENT decides -- \thmref{c02-rewrite-frame}]
    TreeNode* L = dfs(v->left  /*, state*/);   // recurse -> new left root
    TreeNode* R = dfs(v->right /*, state*/);   // recurse -> new right root
    v->left = L;  v->right = R;            // SLOT 2: the rewrite (post-order
    // v->left = R; v->right = L;         //   default) e.g. INVERT = this swap
    return v;                              // SLOT 3: new root of subtree(v);
}                                           //   or nullptr to DELETE, or an
                                            //   orphan pushed to a forest list
TreeNode* transform(TreeNode* root){ return dfs(root); }
\end{lstlisting}

\begin{complexitybox}
Every variant visits each node once: $\bigO{n}$ time, $\bigO{h}$ stack by
\lemref{c02-visit-count}. The rewrite itself is $\bigO{1}$. The only recurring
cost \emph{bug} is an accidental second traversal -- re-flipping to ``fix'' an
orientation instead of fixing SLOT 2 (\thmref{c02-invert-invol} says a double flip
is a no-op), or rebuilding a whole subtree per node when one pass suffices.
\end{complexitybox}

\subsection{Worked Examples}

\begin{example}{Inverting $R$, one call at a time}{c02-invert-trace}
With \code{dfs(v)} $=$ \{ recurse $L$, recurse $R$, \code{swap(v->left,v->right)},
\code{return v} \}: the leaves $4,5,6$ each recurse on two nulls, swap nothing, and
return themselves. Node $2$ gets $L{=}4,R{=}5$, swaps them so its children become
$(5,4)$, returns $2$. Node $3$ gets $L{=}\nil,R{=}6$, swaps to $(6,\nil)$, returns
$3$. The root $1$ gets $L{=}2,R{=}3$, swaps to $(3,2)$, returns $1$. The assembled
mirror matches \figref{c02-invert}.
\end{example}

\begin{tracebox}[Return values up the stack on $R$ (invert)]
\begin{center}\renewcommand{\arraystretch}{1.15}
\begin{tabular}{@{}c c c c c@{}}
\toprule
\textbf{node} & \textbf{children in} & \textbf{$L,R$ returned} & \textbf{after swap} & \textbf{returns}\\
\midrule
$4,5,6$ & leaves   & $(
il,
il)$ / $(
il,
il)$ & unchanged & self\\
$2$     & $4,5$    & $(4,5)$   & $(5,4)$   & $2$\\
$3$     & $(\cdot),6$ & $(\nil,6)$ & $(6,\nil)$ & $3$\\
$1$     & $2,3$    & $(2,3)$   & $(3,2)$   & $1$\\
\bottomrule
\end{tabular}\end{center}
\end{tracebox}

\subsection{Variations and Twists}

\begin{itemize}[nosep]
  \item Flip at \emph{every} node: invert (\figref{c02-invert}); flip at
        \emph{chosen} nodes to hit a target traversal: \secref{c02-transform} lab.  \hfill\lv{3}
  \item Return \code{nullptr} to delete a whole qualifying subtree, decided by its
        aggregate: prune (\probref{c02-prune}) is the template.                    \hfill\lv{3}
  \item Return \code{nullptr} \emph{and} record the orphan as a new tree root:
        delete-and-return-forest (\probref{c02-delete-forest}).                    \hfill\lv{4}
  \item Do not mutate -- \emph{assemble} a tree from descriptions/traversals:
        \probref{c02-create-tree}, construct-from-pre+in (practice).               \hfill\lv{3}
  \item Decide equality \emph{up to} a set of flips -- compare flip-invariant
        multiset signatures: \probref{c02-flip-eq}.                                \hfill\lv{4}
  \item Layout transformation -- the tree is read-only but the \emph{output grid}
        is built by recursion: print in a matrix (\probref{c02-print}).            \hfill\lv{3}
\end{itemize}

\subsection{Common Pitfalls}\label{sec:c02-transform-pitfalls}

\begin{warningbox}[Pitfall 1 -- swapping values instead of pointers]
An ``invert'' that copies $\mathit{val}$ between nodes quietly fails on
non-full trees and any rebuild that must keep structure. A flip is a pointer
exchange (\defref{c02-flip}); touch the $\mathit{val}$ field and you have changed
the problem.
\end{warningbox}
\begin{lstlisting}[caption={\textcolor{warningred}{WRONG}: permutes \texttt{val}s -- happens to look right only on perfect trees}, language=C++, basicstyle=\ttfamily\small]
void inv(TreeNode* r){                        // LC 226-style invert
    if(!r) return;
    std::swap(r->val, (r->left? r->left->val : r->val)); // messes values,
    inv(r->left); inv(r->right);                          // not structure
}                                             // FAILS once children differ in depth
\end{lstlisting}
\begin{lstlisting}[caption={\textcolor{successgreen}{RIGHT}: one pointer swap per node; \thmref{c02-invert-invol} makes it total}, language=C++, basicstyle=\ttfamily\small]
TreeNode* inv(TreeNode* r){
    if(!r) return nullptr;
    TreeNode* L = inv(r->left), *R = inv(r->right);
    r->left = R; r->right = L;                 // swap the POINTERS
    return r;
}
\end{lstlisting}

\begin{warningbox}[Pitfall 2 -- rewriting before recursing (stale children)]
Doing the parent's structural decision using a child's result \emph{before} the
child has been transformed reads the old shape. \thmref{c02-rewrite-frame} is the
rule: child-dependent rewrite goes \emph{after} the calls.
\end{warningbox}
\begin{lstlisting}[caption={\textcolor{warningred}{WRONG}: decides from untransformed children, then recurses}, language=C++, basicstyle=\ttfamily\small]
TreeNode* bad(TreeNode* v){
    if(!v) return nullptr;
    if(prune(v->left)) v->left = nullptr;      // child NOT yet rebuilt here
    v->right = bad(v->right);                  // ordering is inconsistent
    return bad(v->left);                       // double-walks, wrong shape
}
\end{lstlisting}
\begin{lstlisting}[caption={\textcolor{successgreen}{RIGHT}: recurse first, then rewrite from returned roots}, language=C++, basicstyle=\ttfamily\small]
TreeNode* good(TreeNode* v){
    if(!v) return nullptr;
    v->left  = good(v->left);                  // new left root (or null)
    v->right = good(v->right);                 // new right root (or null)
    return keep(v) ? v : nullptr;              // decide on the FINISHED subtree
}
\end{lstlisting}

\begin{warningbox}[Pitfall 3 -- forgetting to reattach the returned pointer]
If you recurse for its side effects but never assign the result back into the
child slot, deletions and rebuilds are silently lost -- the old edges survive.
The contract (\defref{c02-return-newroot}) is \emph{assign then return}.
\end{warningbox}
\begin{lstlisting}[caption={\textcolor{warningred}{WRONG}: calls dfs but discards what it returned}, language=C++, basicstyle=\ttfamily\small]
void leak(TreeNode* v){
    if(!v) return;
    dfs(v->left); dfs(v->right);               // results thrown away:
    if(delete_me(v)) { /* forgot v->parent->x = nullptr; */ }
}
\end{lstlisting}
\begin{lstlisting}[caption={\textcolor{successgreen}{RIGHT}: caller plugs the returned root back into the slot}, language=C++, basicstyle=\ttfamily\small]
TreeNode* dfs(TreeNode* v){
    if(!v) return nullptr;
    v->left  = dfs(v->left);                   // <-- reattach
    v->right = dfs(v->right);                  // <-- reattach
    return delete_me(v) ? nullptr : v;         // parent reattaches THIS
}
\end{lstlisting}


\subsection{Problem Lab (fully solved)}

The ladder runs from a read-only layout build (\probref{c02-print}) and a
from-scratch reconstruction (\probref{c02-create-tree}) up to two genuine pointer
rewrites: equality-up-to-flips (\probref{c02-flip-eq}) and the delete-and-split
summit of the section (\probref{c02-delete-forest}).

\solvedproblem{c02-print}{Print Binary Tree}{LeetCode 655}
  {https://leetcode.com/problems/print-binary-tree/}{LC Medium L3}
  {recursive layout, index arithmetic, output-grid construction}

\begin{problemstatement}[Print Binary Tree]
Given the root of a binary tree, build an $m\times n$ matrix of strings showing a
\emph{formatted layout}. Let $h$ be the tree height (edges root$\to$deepest leaf).
Then $m=h+1$ rows and $n=2^{h+1}-1$ columns. The root sits in the top middle,
\code{res[0][(n-1)/2]}. A node shown at row $r$, column $c$ places its left child
at row $r+1$, column $c-2^{\,h-r-1}$, and its right child at column
$c+2^{\,h-r-1}$. Every cell not occupied by a node holds \code{""}.
Constraints: $1\le n_{\text{nodes}}\le 2000$,
values $\in[0,200]$, $h\le 14$.
\end{problemstatement}

\begin{attemptgate}[Attempt it yourself first -- \lv{3}, give yourself 20 minutes]
A fair first attempt: figure out the row/column of each node by hand on a 3-level
tree, then read off the pattern.
\peek{This is a \emph{layout transformation} -- the tree is read-only but you
\emph{build an output structure} by recursion (\secref{c02-transform}, last twist).}
\nudge{Precompute $h$ once so every cell index is closed-form; the recursion just
writes $\mathit{val}$ at $(r,c)$ and recurses to $(r{+}1,\,c\pm 2^{h-r-1})$.}
\fullhint{Allocate the full \code{""} grid, run one $\bigO{n}$ DFS that fills it,
return the grid. The full build is below.}
\end{attemptgate}

\solstep{Understanding the Problem}
Nothing here rewires the tree; the object being \emph{constructed} is the string
matrix. The whole task is: for each node, know its $(\text{row},\text{col})$ in the
grid. Row is simply its depth. Column follows the midpoint-split rule: the root
takes the middle column and each level halves the horizontal gap, so a node's two
children sit symmetrically $\pm 2^{h-r-1}$ away. Because the offsets depend on the
\emph{global} height, you must know $h$ before placing anything.

\solstep{Key Observations}
\begin{observation}
The grid is a fixed size $m\times n=(h+1)\times(2^{h+1}-1)$ determined by $h$
alone; precompute $h$ with the bottom-up height of \secref{c02-bottomup}, then fill.
\end{observation}
\begin{observation}
Column of a node is not its inorder index -- it is $\text{parent\_col}\mp
2^{h-\text{row}-1}$, so a node's depth \emph{and} its parent's column are the two
pieces of state to carry down (a top-down carry, \secref{c02-topdown}).
\end{observation}
\begin{observation}
$h\le14\Rightarrow n\le 2^{15}-1=32767$ columns and $m\le15$ rows, at most
$\sim4.9\times10^{5}$ cells -- allocating the grid up front with \code{""} is cheap
and lets filling be pure writes with no resizing.
\end{observation}

\solstep{Think It Through With Me}
\begin{thinkbox}
\deadend{Do an inorder traversal and number the columns $0,1,2,\dots$.}{Inorder
index ignores \emph{depth}: two different trees with the same inorder sequence get
different required layouts (the offset halves each level), so inorder columns land
nodes in the wrong cells.}
\deadend{Recompute each node's column from scratch by counting the path of
left/right choices from the root.}{Correct -- the column is a signed sum of powers
of two -- but $O(nh)$ and fiddly; carrying $(r,c)$ down in one DFS collapses it to
$\bigO{n}$ with one subtraction per edge.}
\aha{Carry $(r,c)$ as parameters: place $\mathit{val}$ at $(r,c)$, recurse left to
$(r{+}1,c-2^{h-r-1})$ and right to $(r{+}1,c+2^{h-r-1})$.}
\end{thinkbox}

\solstep{The Approach Ladder}
\begin{center}\renewcommand{\arraystretch}{1.3}
\begin{tabular}{@{}c p{0.235\linewidth} p{0.28\linewidth} l c@{}}
\toprule
\textbf{\#} & \textbf{Approach} & \textbf{Idea} & \textbf{Time} & \textbf{Verdict} \\
\midrule
1 & \lv{2} recompute column per node & signed power-of-two sum from root & $O(nh)$ & \tle \\
2 & \lv{3} carry $(r,c)$ down & one DFS after precomputing $h$ & $O(n)$ & \ac \\
\bottomrule
\end{tabular}\end{center}
Rung 1 is right-but-slow; rung 2 reuses the same formula, memoised in the call
arguments.

\solstep{Solution}
Compute $h$. Allocate \code{res} as $(h+1)$ rows of $(2^{h+1}-1)$ empty strings.
Call \code{place(root, 0, (n-1)/2, h)}: write \code{to\_string(v->val)} at
\code{res[r][c]}, then recurse on the two children at the offset columns. Return
\code{res}. (See \figref{c02-invert} for the tree; the filled grid is in the trace.)

\solstep{Proof of Correctness}
\begin{proofbox}
By induction on depth. The root is placed at $(0,(n-1)/2)$ as required. Assume a
node $v$ is at the correct \code{res[r][c]}; \code{place} writes its left child at
column $c-2^{h-r-1}$ in row $r+1$ and right child at $c+2^{h-r-1}$ -- exactly the
problem's placement rule -- so the child lands correctly. Every non-null node is
visited once and written to its rule-mandated cell; cells never written stay
\code{""}. Hence the matrix equals the specified layout.
\end{proofbox}

\solstep{Implementation}
\begin{lstlisting}[caption={Print Binary Tree (LeetCode 655) -- carry (row, col) down}, label={lst:c02-print}]
#include <bits/stdc++.h>
using namespace std;

struct TreeNode {
    int val;
    TreeNode *left, *right;
    TreeNode(int x = 0) : val(x), left(nullptr), right(nullptr) {}
};

class Solution {
    int height(TreeNode* v) {                    // bottom-up height (\secref{c02-bottomup})
        return v ? 1 + max(height(v->left), height(v->right)) : -1;
    }
    void place(TreeNode* v, int r, int c, int h,
               vector<vector<string>>& res) {
        if (!v) return;
        res[r][c] = to_string(v->val);
        int off = (1 << (h - r - 1));            // 2^(h-r-1)
        place(v->left,  r + 1, c - off, h, res);
        place(v->right, r + 1, c + off, h, res);
    }
public:
    vector<vector<string>> printTree(TreeNode* root) {
        int h = height(root), n = (1 << (h + 1)) - 1;   // 2^(h+1)-1 cols
        vector<vector<string>> res(h + 1, vector<string>(n, ""));
        place(root, 0, (n - 1) / 2, h, res);
        return res;
    }
};
\end{lstlisting}

\solstep{Trace}
\begin{tracebox}[On $R$ -- $h=2$, so $m=3$ rows, $n=2^{3}-1=7$ cols; root at col 3]
\begin{center}\renewcommand{\arraystretch}{1.1}
\begin{tabular}{@{}c|c|c|c|c|c|c|c@{}}
\toprule
row 0 & . & . & . & \textbf{1} & . & . & .\\
row 1 & . & \textbf{2} & . & . & . & \textbf{3} & .\\
row 2 & \textbf{4} & . & \textbf{5} & . & . & . & \textbf{6}\\
\bottomrule
\end{tabular}\end{center}
Root $1@ (0,3)$. Offset at $r{=}0$ is $2^{2-0-1}=2$: so $2@(1,1)$, $3@(1,5)$.
Offset at $r{=}1$ is $2^{0}=1$: $4@(2,0)$, $5@(2,2)$, and $3$'s lone right child
$6@(2,6)$ (its left slot stays empty). Dots are \code{""}. \ac
\end{tracebox}

\solstep{Complexity}
\begin{complexitybox}
Time $\bigO{n}$ to fill ($n_{\text{nodes}}$ writes) plus $\bigO{n}$ for the height
pass; space $\bigO{m\cdot 2^{h}}$ for the grid -- at $h\le14$ that is $\le 4.9\times10^{5}$
cells, well inside memory. Stack $\bigO{h}$, and $h\le14$ here so no recursion-depth
risk (unlike the general skewed trees later).
\end{complexitybox}

\solstep{Verification and Edge Cases}
\begin{center}\renewcommand{\arraystretch}{1.15}
\begin{tabular}{@{}p{0.34\linewidth} p{0.60\linewidth}@{}}
\toprule
\textbf{Edge case} & \textbf{What the code does} \\
\midrule
Single node & $h=0$, $n=1$: \code{res[0][0]} = value, offset shift $2^{-1}$ never
used (no children). \\
Skewed chain & still fills correctly; only the used cells get strings. \\
Node value $0$ & \code{to\_string(0)="0"} is truthy-as-text -- never confuse the
empty cell \code{""} with \code{"0"}. \\
\bottomrule
\end{tabular}\end{center}

\solstep{Pitfalls and Variants}
\begin{warningbox}
Use \code{int off = 1 << (h-r-1)} \emph{with parentheses}: \code{1 << h - 1} parses
as \code{1 << (h-1)} (shift binds looser than \code{-}) -- a silent off-by-one power
of two. Also keep \code{""} (empty string) distinct from a value's string; a node
of value $0$ must print \code{"0"}, not stay blank.
\end{warningbox}

\begin{insightbox}[Key takeaways -- three things to remember]
\begin{enumerate}[nosep, leftmargin=1.6em]
  \item When the output is a \emph{structure} you build by recursion, carry the
        target coordinates down -- memoise placement in the arguments.
  \item Precompute any \emph{global} quantity (here $h$) the index formula needs;
        one cheap pass up front keeps the filling DFS pure writes.
  \item ``Read the tree, write something else'' is a transformation of the
        \emph{output}, not the pointers -- the same (r,c) carry is a top-down
        (\secref{c02-topdown}) pattern wearing a construction hat.
\end{enumerate}
\end{insightbox}

% ---------------------------------------------------------------------
\solvedproblem{c02-create-tree}{Create Binary Tree From Descriptions}{LeetCode 2196}
  {https://leetcode.com/problems/create-binary-tree-from-descriptions/}{LC Medium L3}
  {reconstruction, hash-indexed nodes, root-by-in-degree}

\begin{problemstatement}[Create Binary Tree From Descriptions]
You are given \code{descriptions}, a list of triples $[\,\mathit{parent},\mathit{child},
\mathit{isLeft}\,]$ meaning ``\code{child} is the left ($\mathit{isLeft}=1$) or right
($0$) child of \code{parent}.'' All node values are distinct and the triples define
exactly one valid binary tree. Build it and return its root. $1\le|\text{descriptions}|\le10^{4}$,
values $\in[1,10^{5}]$.
\end{problemstatement}

\begin{attemptgate}[Attempt it yourself first -- \lv{3}, give yourself 20 minutes]
A fair first attempt: ignore the tree for a moment and ask how to \emph{find the
root} from the triples alone.
\peek{This is a rebuild (\defref{c02-mutate-rebuild}) -- the input is a
\emph{description}, so you allocate fresh nodes rather than editing any given tree.}
\nudge{Every non-root node appears in the \code{child} slot exactly once; the root
never appears as a child. How do you find the one node absent from the child set?}
\fullhint{Index nodes by value in a hashmap (create-on-demand), attach children,
record which values are children, and the value present as a parent but absent
from the child set is the root. Build below.}
\end{attemptgate}

\solstep{Understanding the Problem}
The triples arrive in arbitrary order, so you cannot scan once and assume a parent
exists yet -- you must be able to fetch-or-create a node by value. A valid binary
tree has exactly one node of in-degree $0$: the root. Everything else is decided
once you (a) materialise nodes lazily by value and (b) know which value is the
unique non-child.

\solstep{Key Observations}
\begin{observation}
A \code{unordered\_map<int,TreeNode*>} with ``create if absent'' makes the order of
processing triples irrelevant: parent and child are each resolved through the map,
so the same node object is reused wherever its value reappears.
\end{observation}
\begin{observation}
The root is the only value that appears as a $\mathit{parent}$ but never as a
$\mathit{child}$; tracking the set of child-values finds it in one pass over the map.
\end{observation}

\solstep{Think It Through With Me}
\begin{thinkbox}
\deadend{Treat the triples like an edge list and root a DFS from $\mathit{descriptions}[0][0]$.}{The
first triple's parent need not be the tree's root -- it could itself be some other
node's child -- so you would return a non-root node as the answer.}
\deadend{Sort the triples by value and build top-down.}{Sorting gains nothing: the
root is defined by \emph{not being anyone's child}, a set-membership fact, not an
ordering one.}
\aha{Realise the root is exactly the node with in-degree $0$; a child-set plus the
key set of created nodes exposes it directly.}
\end{thinkbox}

\solstep{The Approach Ladder}
\begin{center}\renewcommand{\arraystretch}{1.3}
\begin{tabular}{@{}c p{0.235\linewidth} p{0.28\linewidth} l c@{}}
\toprule
\textbf{\#} & \textbf{Approach} & \textbf{Idea} & \textbf{Time} & \textbf{Verdict} \\
\midrule
1 & \lv{2} test each node as root & validate a full build per candidate & $O(n^2)$ & \tle \\
2 & \lv{3} child-set + hashmap & root = the value never a child & $O(n)$ & \ac \\
\bottomrule
\end{tabular}\end{center}
Rung 1 rebuilds the tree for every suspected root; rung 2 reads the answer off the
in-degree-$0$ node in one pass.

\solstep{Solution}
Maintain \code{node} (value $\to$ pointer, created on demand) and \code{isChild}
(set of child values). For each triple resolve parent and child through
\code{node}, insert the child value into \code{isChild}, and wire
\code{p->left}/{\code{->right}$=c$}. Finally return \code{node[r]} where $r$ is the
unique key not in \code{isChild}.

\solstep{Proof of Correctness}
\begin{proofbox}
Create-on-demand makes \code{node[x]} the \emph{same} object for every occurrence of
value $x$, so each triple attaches the correct existing nodes; distinct values give
each node at most one parent, and the input is promised valid, so no node is
overwritten. In a tree every node except the root is a child of exactly one parent,
so the child set contains all values but the root. The keys of \code{node} are all
values; their set difference is a singleton $\{r\}$, and $r$ is the root by
in-degree $0$. Returning \code{node[r]} is therefore correct.
\end{proofbox}

\solstep{Implementation}
\begin{lstlisting}[caption={Create Binary Tree From Descriptions (LeetCode 2196) -- lazy node map + child set}, label={lst:c02-create-tree}]
#include <bits/stdc++.h>
using namespace std;

struct TreeNode {
    int val;
    TreeNode *left, *right;
    TreeNode(int x = 0) : val(x), left(nullptr), right(nullptr) {}
};

class Solution {
public:
    TreeNode* createBinaryTree(vector<vector<int>>& desc) {
        unordered_map<int, TreeNode*> node;
        unordered_set<int> isChild;
        auto get = [&](int v) {                    // fetch-or-create by value
            if (!node.count(v)) node[v] = new TreeNode(v);
            return node[v];
        };
        for (auto& d : desc) {
            TreeNode* p = get(d[0]);
            TreeNode* c = get(d[1]);
            isChild.insert(d[1]);
            if (d[2]) p->left = c; else p->right = c;
        }
        for (auto& kv : node)
            if (!isChild.count(kv.first)) return kv.second;   // in-degree 0
        return nullptr;
    }
};
\end{lstlisting}

\solstep{Trace}
\begin{tracebox}[$R$ as descriptions (order shuffled): the root is the value absent from the child set]
Triples: $[2,5,0],[1,2,1],[3,6,0],[1,3,0],[2,4,1]$. Child set $=\{2,5,6,3,4\}$.
Created values $=\{1,2,3,4,5,6\}$. Only $1$ is a parent but never a child, so the
map lookup for $1$ is returned -- and it carries $2$ left, $3$ right, matching
$R$'s root. \ac
\end{tracebox}

\solstep{Complexity}
\begin{complexitybox}
Time $\bigO{n}$ expected (hashmap/set operations, $n=|\text{descriptions}|$),
space $\bigO{n}$ for the nodes plus the child set. No recursion, so stack is
$\bigO{1}$ regardless of shape -- the safe default when the description could
describe a deep chain.
\end{complexitybox}

\solstep{Verification and Edge Cases}
\begin{center}\renewcommand{\arraystretch}{1.15}
\begin{tabular}{@{}p{0.34\linewidth} p{0.60\linewidth}@{}}
\toprule
\textbf{Edge case} & \textbf{What the code does} \\
\midrule
Root appears only as a parent & never enters \code{isChild}, so it is the answer. \\
Triples out of order & create-on-demand resolves parent/child regardless. \\
Two-leaf, single-child mixes & distinct values forbid a node being its own
ancestor; input is promised valid. \\
\bottomrule
\end{tabular}\end{center}

\solstep{Pitfalls and Variants}
\begin{warningbox}
Do not assume \code{node[]} is fully populated when you first see a parent -- that
is exactly the bug create-on-demand cures. If the problem later asks you to
\emph{transform} the built tree (invert, prune), you now have a real root and can
run the \secref{c02-transform} skeleton over it; construction and transformation
compose cleanly.
\end{warningbox}

\begin{insightbox}[Key takeaways -- three things to remember]
\begin{enumerate}[nosep, leftmargin=1.6em]
  \item Rebuild by \emph{indexing nodes by value} with fetch-or-create; it erases
        any dependence on the input's order.
  \item ``Find the root'' is almost always ``find the in-degree-$0$ node'' -- one
        set of children answers it.
  \item Construction returns a root pointer just like a transformation does; the
        recursion discipline (\defref{c02-return-newroot}) is shared even when the
        build itself is iterative.
\end{enumerate}
\end{insightbox}


% ---------------------------------------------------------------------
\solvedproblem{c02-flip-eq}{Flip Equivalent Binary Trees}{LeetCode 951}
  {https://leetcode.com/problems/flip-equivalent-binary-trees/}{LC Medium L4}
  {structural equivalence, two-case recursion, value-uniqueness shortcut}

\begin{problemstatement}[Flip Equivalent Binary Trees]
A \emph{flip operation} picks one node and swaps its left and right child
subtrees. Two trees are \emph{flip equivalent} if one can be turned into the other
by a sequence of such flips. Given roots $a$ and $b$, return \code{true} iff they
are flip equivalent. Node values are unique within each tree; $0\le$\,nodes$\le100$.
\end{problemstatement}

\begin{attemptgate}[Attempt it yourself first -- \lv{4}, give yourself 25 minutes]
A fair first attempt: at a single node, ask what \code{flipEquiv} must check about
the two children -- there are only two ways to pair them.
\peek{This is an equivalence-under-flip test (\secref{c02-transform}, fifth twist):
a shape question, but you build nothing -- you \emph{compare}.}
\nudge{At matched nodes, either the children correspond straight or they
correspond crossed (one flip). Write the recursion over those two pairings; the
answer is the disjunction.}
\fullhint{\code{flipEquiv}(a,b) is: values equal \textbf{and} either the two
children match straight or they match crossed (one flip here). Uniqueness of
values makes the short-circuit $\bigO{n}$. Build below.}
\end{attemptgate}

\solstep{Understanding the Problem}
Flips are independent per node and can be applied in any order (\thmref{c02-invert-invol}
style argument), so the only question is local: when two nodes match in value, can
we pair their children -- either left-with-left and right-with-right (no flip
here), or left-with-right and right-with-left (a flip here)? That binary choice at
every node is the whole problem. Uniqueness of values means at most one of the two
pairings can ever succeed, which is what keeps the recursion linear.

\solstep{Key Observations}
\begin{observation}
$a$ and $b$ are flip equivalent iff (both null) or (one null, values equal, and
their children match either straight or crossed). Base: two nulls are equivalent,
exactly one null is not.
\end{observation}
\begin{observation}
Because values are unique, if the straight pairing holds the crossed pairing
cannot (the child roots would carry the same two values in both orders only if
they were equal), so short-circuiting the \code{||} visits each node $O(1)$ times.
\end{observation}
\begin{observation}
No tree is mutated or built -- this decides \emph{equality up to flips} by
comparing a flip-invariant (the multiset of child-subtree identities), matching the
pattern card's third use (\defref{c02-flip}).
\end{observation}

\solstep{Think It Through With Me}
\begin{thinkbox}
\deadend{Canonicalise each tree by always flipping so the smaller-valued child is
on the left, then compare structurally.}{Works, but it is two rebuilds plus a
comparison -- more code and an $\bigO{n}$ allocation per tree, when a single
in-place recursion decides the same thing with no extra memory.}
\deadend{Compare heights and sizes to reject quickly.}{Necessary-not-sufficient:
many non-equivalent trees share the height/size profile, so it only filters; it
never proves equivalence.}
\aha{The equivalence is defined \emph{recursively}, so test it recursively -- at
each matched node just try the two child pairings and \code{or} them.}
\end{thinkbox}

\solstep{The Approach Ladder}
\begin{center}\renewcommand{\arraystretch}{1.3}
\begin{tabular}{@{}c p{0.235\linewidth} p{0.28\linewidth} l c@{}}
\toprule
\textbf{\#} & \textbf{Approach} & \textbf{Idea} & \textbf{Time} & \textbf{Verdict} \\
\midrule
1 & \lv{3} enumerate flip subsets & try all $2^{n}$ flip choices & $O(2^{n}\!\cdot\!n)$ & \tle \\
2 & \lv{4} two-case recursion & straight \textbf{or} crossed at each node & $O(n)$ & \ac \\
\bottomrule
\end{tabular}\end{center}
Rung 1 is the definition unrolled; rung 2 factors it -- the choice at a node is
independent of the rest, so we branch once, not per subset.

\solstep{Solution}
Structural recursion \code{fe(a,b)}: return \code{true} if both nodes are null,
\code{false} if exactly one is null or the values differ; otherwise \code{fe} is
the disjunction of two pairings -- the \emph{straight} one, joining
\code{fe(a->left, b->left)} with \code{fe(a->right, b->right)}, and the
\emph{crossed} one, joining \code{fe(a->left, b->right)} with
\code{fe(a->right, b->left)}. Short-circuit keeps it linear (Observation 2).

\solstep{Proof of Correctness}
\begin{proofbox}
Induction on $\size(a)+\size(b)$. Base cases (both/null/value-mismatch) are
immediate. Step, matched values: if a flip sequence $F$ makes $a\equiv b$, then at
the roots either the root is not flipped by $F$ (so $a$'s left $\equiv$ $b$'s left
and rights correspond) or it is flipped once at the root (crossed correspondence);
in either case the relevant pairing holds and, by the inductive hypothesis, the
matching recursive calls return \code{true}, so \code{fe} returns \code{true}.
Conversely, if \code{fe} returns \code{true} via the straight pairing, apply the
inductively-guaranteed flips inside each child subtree -- $a\to b$ with no root
flip; via the crossed pairing, additionally flip the root -- a valid flip sequence.
Hence \code{fe} is correct.
\end{proofbox}

\solstep{Implementation}
\begin{lstlisting}[caption={Flip Equivalent Binary Trees (LeetCode 951) -- straight-or-crossed recursion}, label={lst:c02-flip-eq}]
#include <bits/stdc++.h>
using namespace std;

struct TreeNode {
    int val;
    TreeNode *left, *right;
    TreeNode(int x = 0) : val(x), left(nullptr), right(nullptr) {}
};

class Solution {
public:
    bool flipEquiv(TreeNode* a, TreeNode* b) {
        if (!a && !b) return true;                       // both empty
        if (!a || !b || a->val != b->val) return false;  // shape/value mismatch
        bool straight = flipEquiv(a->left,  b->left)
                     && flipEquiv(a->right, b->right);
        bool crossed  = flipEquiv(a->left,  b->right)
                     && flipEquiv(a->right, b->left);    // one flip at this node
        return straight || crossed;                      // short-circuits: O(n)
    }
};
\end{lstlisting}

\solstep{Trace}
\begin{tracebox}[$R$ vs.\ $S$, where $S$ is $R$ with a single flip at node $2$ (its children become $5,4$)]
Root: $\mathit{fe}(1,1)$ values equal. Straight pairing recurses
$\mathit{fe}(2,2)$ and $\mathit{fe}(3,3)$. Node $3$ matches (both have lone right
$6$) $\to$ \code{true}. At $\mathit{fe}(2,2)$: children are $(4,5)$ in $R$ and
$(5,4)$ in $S$. Straight: $\mathit{fe}(4,5)$ = false (values differ). Crossed:
$\mathit{fe}(4,4)\,\&\&\,\mathit{fe}(5,5)$ = true. So node $2$ returns true, hence
the root's straight branch is true, hence the trees are flip equivalent -- one flip
at node $2$ suffices. \ac
\end{tracebox}

\solstep{Complexity}
\begin{complexitybox}
Time $\bigO{n}$: uniqueness (Observation 2) means only one of \{straight,
crossed\} can match, so the \code{||}/\code{\&\&} short-circuit prevents the naive
$4$-branching; each node is visited a constant number of times. Space $\bigO{h}$
stack. With $n\le100$ this is never near the limits, but the \emph{reason} it is
linear (uniqueness) must be re-checked if values were allowed to repeat -- then
both pairings may match and you get $O(2^{h})$ without memoisation on subtree
identities.
\end{complexitybox}

\solstep{Verification and Edge Cases}
\begin{center}\renewcommand{\arraystretch}{1.15}
\begin{tabular}{@{}p{0.34\linewidth} p{0.60\linewidth}@{}}
\toprule
\textbf{Edge case} & \textbf{What the code does} \\
\midrule
Both empty & returns \code{true} at the base. \\
One empty, one not & returns \code{false} (asymmetric). \\
Identical trees & straight branch true everywhere, no flip needed. \\
Perfect mirror & crossed branch fires at every node; still \emph{true}. \\
Repeated values (not in constraints) & would need memo/canonical form -- noted
in the complexity box. \\
\bottomrule
\end{tabular}\end{center}

\solstep{Pitfalls and Variants}
\begin{warningbox}
Comparing only the child \emph{values} at each node (does $\{a_L,a_R\}=\{b_L,b_R\}$
as a set) is \emph{not} enough -- the subtrees under matched children must
themselves be flip equivalent; recurse into them. Variant: ``minimum flips to make
equivalent'' replaces the \code{||} with a $\min$ over the two pairings, adding $1$
on the crossed branch -- same skeleton, counting instead of deciding.
\end{warningbox}

\begin{insightbox}[Key takeaways -- three things to remember]
\begin{enumerate}[nosep, leftmargin=1.6em]
  \item ``Equal up to independent per-node choices'' problems collapse to a
        recursion that tries each local choice and \code{or}/\code{min}s them.
  \item A side-condition (value uniqueness) is often what turns an exponential
        branch-and-bound into a linear scan -- name it, do not assume it.
  \item Equivalence checks are transformations' cousins: the flip is the same
        pointer swap (\defref{c02-flip}), here only imagined, never applied.
\end{enumerate}
\end{insightbox}

\begin{hardeningbox}
\begin{itemize}[nosep, leftmargin=1.4em]
  \item If values are \emph{not} unique, memoise $\mathit{fe}$ on
        $(\mathit{ptr}(a),\mathit{ptr}(b))$ (or canonical hash of each subtree) to
        kill the exponential re-comparison of matched shapes.
  \item Deep skewed input ($n$ up to $10^{5}$ variants exist) $\to$ the two-sided
        recursion is $\bigO{h}$ stack: convert to an explicit pair-stack
        $\{(a,b)\}$ worklist, pushing the two candidate pairings.
\end{itemize}
\end{hardeningbox}

% ---------------------------------------------------------------------
\solvedproblem{c02-delete-forest}{Delete Nodes And Return Forest}{LeetCode 1110}
  {https://leetcode.com/problems/delete-nodes-and-return-forest/}{LC Medium L4}
  {post-order rewire, return-new-root, orphan recording, set membership}

\begin{problemstatement}[Delete Nodes And Return Forest]
Given the root of a binary tree and a set of distinct values \code{toDelete},
delete every node whose value is in the set (removing it and its parent edge). The
children of a deleted node, if they survive, become roots of new trees. Return the
list of roots of the resulting forest, in any order. $1\le$\,nodes$\le1000$,
values in $[1,10000]$ distinct, $|\text{toDelete}|\le100$.
\end{problemstatement}

\begin{attemptgate}[Attempt it yourself first -- \lv{4}, give yourself 25 minutes]
A fair first attempt: decide \emph{when} you know a node must be reported as a new
forest root -- it is a child of a deleted node that was not itself deleted.
\peek{This is the return-the-new-root contract (\defref{c02-return-newroot}) plus
Pitfall 3: deletion is ``return \code{nullptr} to the parent''; a split is ``also
record this surviving child as a root.''}
\nudge{Process children \emph{first} (post-order, \thmref{c02-rewrite-frame}) so a
deleted node's children already point at their rebuilt subtrees; then its
non-null children are exactly the new roots.}
\fullhint{\code{dfs(v)} reattaches \code{v->left=dfs(left)}, \code{v->right=dfs(right)};
if $v$ is deleted, push its surviving children to the answer and return
\code{nullptr}; else return $v$. After the call, push the root if it survived. Below.}
\end{attemptgate}

\solstep{Understanding the Problem}
Deleting a node changes the pointer graph, so this is a genuine transformation --
and it produces \emph{several} outputs (a forest), not one tree. Two facts drive
the design: (i) a node's fate must be decided \emph{after} its children are
processed, because a deleted node hands its (already-rebuilt) surviving children
to the forest; (ii) the parent must be told the child vanished, which is exactly
what ``return \code{nullptr} from the child call and assign it into the slot''
does. Membership in \code{toDelete} should be $O(1)$ -- hash the set.

\solstep{Key Observations}
\begin{observation}
A node becomes a forest root iff it survives and its parent is deleted (or it is
the overall root and survives). So the answer $=$ survivors whose parent edge is
cut, collected as we go.
\end{observation}
\begin{observation}
Post-order is forced: only after \code{dfs} on both children do we know their
final roots; if $v$ is then deleted, its current non-null children \emph{are} the
new roots (their pointers already exclude deletions below them).
\end{observation}
\begin{observation}
Reattaching the returned pointer (\code{v->left = dfs(v->left)}) is what actually
severs a deleted grandchild -- forgetting it leaves dangling edges (Pitfall 3).
\end{observation}

\solstep{Think It Through With Me}
\begin{thinkbox}
\deadend{Pre-order: if $v$ is deleted, splice it out, then recurse.}{When you cut
$v$ you do not yet know which of $v$'s descendants will also be deleted, so you
record the \emph{wrong} children as forest roots (a child that is itself deleted
would be reported). \thmref{c02-rewrite-frame} forbids a child-dependent rewrite
before the children run.}
\deadend{For each value in \code{toDelete}, do a full DFS to find and remove it.}{$O(n\cdot|\text{toDelete}|)$
and awkward: after removing one node the parent pointers you relied on may already
be gone; interleaved deletions corrupt the walk. One post-order pass handles all
deletions at once in $\bigO{n}$.}
\aha{Return the rebuilt subtree root; \code{nullptr} means ``I was deleted'', and a
deleted node's surviving children are pushed to the forest on the way out.}
\end{thinkbox}

\solstep{The Approach Ladder}
\begin{center}\renewcommand{\arraystretch}{1.3}
\begin{tabular}{@{}c p{0.235\linewidth} p{0.28\linewidth} l c@{}}
\toprule
\textbf{\#} & \textbf{Approach} & \textbf{Idea} & \textbf{Time} & \textbf{Verdict} \\
\midrule
1 & \lv{3} per-target full DFS & one pass to delete each value & $O(n\,k)$ & \tle \\
2 & \lv{3} pre-order splice & cut then recurse & $O(n)$ & \textcolor{warningred}{\ding{55}\ \textbf{WA}} \\
3 & \lv{4} post-order return-null & rebuild children, then decide & $O(n)$ & \ac \\
\bottomrule
\end{tabular}\end{center}
Rung 2 is the trap \thmref{c02-rewrite-frame} warns about (wrong forest roots);
rung 3 fixes it by deciding after the children.

\solstep{Solution}
Hash \code{toDelete} into a set. Post-order \code{dfs(v)}: recurse and \emph{assign}
both children; if $v$'s value is deleted, push each surviving (non-null) child to
\code{forest} and return \code{nullptr}; else return $v$. Call it on the root and,
if the returned root is non-null (root survived), push it. Return \code{forest}.

\solstep{Proof of Correctness}
\begin{proofbox}
Induction on $\size$ (post-order, \thmref{c02-postorder}). For $v=\nil$ return
\code{nullptr}, correct. For internal $v$: by hypothesis \code{dfs} returns the
root of the correctly-forested $\subtree(L)$ and $\subtree(R)$; assigning them in
keeps $v$ attached to exactly the survivors that remain in its subtrees. If $v$ is
\emph{not} deleted, $v$ is the root of the rebuilt $\subtree(v)$ -- return $v$
(correct). If $v$ \emph{is} deleted, $\subtree(v)$ splits into the two already-rebuilt
child trees; a non-null child root is by the inductive hypothesis the root of a
maximal surviving subtree, so pushing precisely those (and returning
\code{nullptr} to detach $v$ from its parent) yields the correct forest fragment.
The top-level push adds the root's tree iff the root survived. Every surviving
node is in exactly one reported tree (its nearest surviving ancestor whose parent
edge was cut is the reported root), and no deleted node is a root -- so the
collection is exactly the forest.
\end{proofbox}

\solstep{Implementation}
\begin{lstlisting}[caption={Delete Nodes And Return Forest (LeetCode 1110) -- post-order, return-new-root}, label={lst:c02-delete-forest}]
#include <bits/stdc++.h>
using namespace std;

struct TreeNode {
    int val;
    TreeNode *left, *right;
    TreeNode(int x = 0) : val(x), left(nullptr), right(nullptr) {}
};

class Solution {
    vector<TreeNode*> forest;
    unordered_set<int> del;
    TreeNode* dfs(TreeNode* v) {
        if (!v) return nullptr;
        v->left  = dfs(v->left);            // rebuild children FIRST (post-order)
        v->right = dfs(v->right);           // reattach -> severed edges vanish
        if (del.count(v->val)) {            // delete v: its survivors split off
            if (v->left)  forest.push_back(v->left);
            if (v->right) forest.push_back(v->right);
            return nullptr;                 // tell parent to detach v
        }
        return v;                           // v survives as subtree root
    }
public:
    vector<TreeNode*> delNodes(TreeNode* root, vector<int>& toDelete) {
        del = unordered_set<int>(toDelete.begin(), toDelete.end());
        if (dfs(root)) forest.push_back(root);   // root survived
        return forest;
    }
};
\end{lstlisting}

\solstep{Trace}
\begin{tracebox}[$R$ with toDelete containing $3$ (the highlighted chain $1\to3\to6$ in \figref{c02-running})]
$\mathit{dfs}(1)$: left $=\mathit{dfs}(2)$ returns $2$ (children $4,5$ intact, $2$
kept); right $=\mathit{dfs}(3)$: node $3$'s left is null, its right $=\mathit{dfs}(6)$
$=6$; $3$ is deleted $\to$ push surviving child $6$ to forest, return \code{nullptr}.
Back at $1$: \code{v->right=nullptr} severs the edge; $1$ kept, return $1$. The
top-level call sees a non-null root, pushes $1$. Forest $=\{6,\,1\}$ -- the tree
rooted at $1$ (now right-childless) and the orphaned singleton $6$. \ac
\end{tracebox}

\solstep{Complexity}
\begin{complexitybox}
Time $\bigO{n}$ (one post-order; each \code{del.count} is $O(1)$ expected), space
$\bigO{n+|\text{toDelete}|}$ for the forest vector and hash set plus $\bigO{h}$
stack. At $n\le1000$ no depth risk; a $10^{5}$-node chain would need the iterative
form (hardening box).
\end{complexitybox}

\solstep{Verification and Edge Cases}
\begin{center}\renewcommand{\arraystretch}{1.15}
\begin{tabular}{@{}p{0.34\linewidth} p{0.60\linewidth}@{}}
\toprule
\textbf{Edge case} & \textbf{What the code does} \\
\midrule
Root deleted & top-level \code{dfs} returns \code{nullptr}; only its surviving
children were pushed -- root not in forest. \\
Delete a whole branch & its surviving leaves under it are pushed at the highest
deleted node exactly once. \\
Nothing deleted & every call returns its node; forest $=\{$root$\}$. \\
Adjacently-deleted parent/child & child returns \code{nullptr} to parent first
(post-order), so a deleted child is \emph{not} pushed by its deleted parent. \\
\bottomrule
\end{tabular}\end{center}

\solstep{Pitfalls and Variants}
\begin{warningbox}
Reporting forest roots in \emph{pre-order} (before children) double-counts: a node
whose parent is deleted is a root \emph{only if it survives}, which you learn after
processing it. Test adjacent deletions $\{p, \mathit{child}(p)\}$ -- a correct
post-order push nothing for the deleted child, while a naive version pushes it.
Variant: ``return the forest as a sorted list of root values'' just maps the answer
through $\mathit{val}$; ``delete leaves only'' guards the push on childlessness.
\end{warningbox}

\begin{insightbox}[Key takeaways -- three things to remember]
\begin{enumerate}[nosep, leftmargin=1.6em]
  \item A transformation that \emph{splits} the tree is still ``return the new
        subtree root''; extra outputs are collected in a side list as you cut.
  \item Post-order is not a style choice here -- child-dependent structural
        decisions are only safe once the children are final (\thmref{c02-rewrite-frame}).
  \item \code{nullptr}-as-signal is the whole deletion trick: the parent's
        \code{child = dfs(child)} assignment does the surgery for free.
\end{enumerate}
\end{insightbox}

\begin{hardeningbox}
\begin{itemize}[nosep, leftmargin=1.4em]
  \item $n$ grows to $10^{5}$ as a chain: $\bigO{h}$ recursion \textbf{RE}s. Use an
        explicit-stack post-order (push $\{v,\text{visited}\}$, and on the second
        visit assign children and apply the delete rule) so the ``children first''
        order survives without the call stack.
  \item If the same node objects must be \emph{reused} elsewhere, \code{delete} the
        removed nodes (LeetCode hands you ownership, so it is optional there, but
        not in your own code) -- or leaks accumulate exactly where the edges were
        cut.
\end{itemize}
\end{hardeningbox}


\subsection{Practice Problems}

\begin{stuckbox}
\begin{itemize}[nosep, leftmargin=1.4em]
  \item Does your recursion \emph{return the (new) subtree root} and does the
        caller \emph{assign} it back into the child slot? (If not, deletions and
        rebuilds are silently lost -- Pitfall 3, \defref{c02-return-newroot}.)
  \item Is your per-node rewrite placed on the right side of the recursion --
        after the calls when it reads children, before them only when the parent
        decides? (\thmref{c02-rewrite-frame})
  \item Are you swapping \emph{pointers} or \emph{values}? (Invert must touch
        pointers; \defref{c02-flip}.)
  \item For a rebuild: have you found the root by in-degree $0$, and can you
        fetch-or-create any node by value regardless of input order?
        (\probref{c02-create-tree})
  \item If you compare shapes under flips, did you recurse \emph{into} the
        matched subtrees rather than just checking child values? (\probref{c02-flip-eq})
\end{itemize}
\end{stuckbox}

\begin{practiceset}
  \practice{Flip Binary Tree To Match Preorder Traversal}{LeetCode 971}
    {https://leetcode.com/problems/flip-binary-tree-to-match-preorder-traversal/}{LC Medium L3}
    {preorder walk where at each node you may swap children to match the target
    sequence; count the flips, return $-1$ on impossibility -- the
    parent-decides/pre-order side of \thmref{c02-rewrite-frame} with a global flip
    counter (\probref{c02-flip-eq}'s mirror image).}
  \practice{Maximum Binary Tree}{LeetCode 654}
    {https://leetcode.com/problems/maximum-binary-tree/}{LC Medium L3}
    {rebuild: the max of a range is the subtree root, left subarray builds the left
    subtree, right subarray the right -- a construction returning the new root
    (\defref{c02-mutate-rebuild}); monotonic-stack $\bigO{n}$ is the stretch.}
  \practicehh{Construct Binary Tree from Preorder and Inorder Traversal}{LeetCode 105}
    {https://leetcode.com/problems/construct-binary-tree-from-preorder-and-inorder-traversal/}{LC Medium L4}
    {preorder gives the root, inorder splits into left/right extents; recurse with
    index ranges and return the assembled root.}
    {Hash the inorder values to indices for $O(1)$ splits; carry a moving preorder
    pointer.}
    {Guard the empty-range base case; the left subtree size is
    $\mathit{inIdx}-\mathit{inLeft}$ -- compute it before recursing (compare
    \probref{c02-create-tree}'s root-by-in-degree).}
\end{practiceset}

\begin{recallbox}[10-second refresher -- 2.4 Structural transformation]
Point the recursion at the tree: each call \emph{returns the new root} of the
subtree it rebuilt (a node, or \code{nullptr} to delete, or a node plus an orphan
pushed to a forest). Child-dependent rewrites go post-order, parent-decided ones
pre-order (\thmref{c02-rewrite-frame}); flip $\ne$ move-values (\thmref{c02-invert-invol}).
\selftest{
  \item State the return-the-new-root contract and name the three things a call
        can return (kept root, \code{nullptr}, root-plus-recorded-orphan).
  \item Why must invert swap \emph{pointers}, not \emph{values}, and why is a
        double invert the identity? (\thmref{c02-invert-invol})
  \item Give one problem where the rewrite goes pre-order and one where it must go
        post-order, and cite the deciding rule (\thmref{c02-rewrite-frame}).
  \item In \probref{c02-delete-forest}, why does a pre-order report double-count a
        node whose parent is also deleted?
  \item Write the transform skeleton (\coderef{c02-transform-skeleton}) from memory,
        marking SLOT 2 (rewrite) and SLOT 3 (return).
}
\end{recallbox}
